\title{Controlling Text-to-Image Diffusion by Orthogonal Finetuning}

\begin{abstract}
Large-scale text-to-image diffusion models demonstrate remarkable ability to produce photorealistic images from textual descriptions. A key open challenge is how to effectively steer or control these powerful models to accomplish various downstream tasks. To address this, we propose a principled finetuning approach called Orthogonal Finetuning~(OFT), designed to adapt text-to-image diffusion models for downstream applications. Unlike prior methods, OFT provably preserves hyperspherical energy, which characterizes the pairwise relationships among neurons on the unit hypersphere. We identify this preservation as critical to maintaining the semantic generation capabilities of text-to-image diffusion models. To enhance finetuning stability, we introduce Constrained Orthogonal Finetuning (COFT), which adds a radius constraint to the hypersphere. In particular, we focus on two key finetuning tasks: subject-driven generation, aiming to create images of a specific subject in new contexts given a few images and a text prompt; and controllable generation, which enables the model to incorporate additional control signals. Our empirical results demonstrate that OFT surpasses existing methods in both generation quality and convergence speed.
\end{abstract}

\section{Introduction}

Recent advancements in text-to-image diffusion models~\cite{saharia2022photorealistic,ramesh2022hierarchical,rombach2022high} have achieved impressive results in producing high-quality images guided by text. However, text prompts can remain ambiguous or insufficient for providing detailed and precise control over the generated images. To overcome these limitations, this paper focuses on two specific types of text-to-image generation tasks:

\begin{itemize}[leftmargin=*,nosep]

    \item \textbf{Subject-driven generation}~\cite{ruiz2023dreambooth}: Given only a few images of a particular subject, the goal is to generate new images of the same subject in different settings using text guidance.
    \item \textbf{Controllable generation}~\cite{zhang2023adding,mou2023t2i}: Provided with an additional control input (\emph{e.g.}, canny edges, segmentation maps), the objective is to produce images that conform to this control signal alongside the text prompt.
\end{itemize}

Both tasks fundamentally revolve around how to finetune text-to-image diffusion models effectively without compromising the generative capabilities acquired during pretraining. We outline the desired qualities for an effective finetuning approach as follows: (1) \emph{training efficiency}: involving fewer trainable parameters and reduced training epochs, and (2) \emph{generalizability preservation}: maintaining high-quality, diverse generative outputs. To achieve this, finetuning is commonly performed either by adjusting neuron weights with a small learning rate (\emph{e.g.}, \cite{ruiz2023dreambooth}) or by incorporating a small module with re-parameterized neuron weights (\emph{e.g.}, \cite{hulora2022,zhang2023adding}). Despite their simplicity, neither approach guarantees that the generative performance from pretraining is retained. Moreover, there is a lack of systematic insight into devising an optimal finetuning strategy and selecting appropriate hyperparameters like the number of epochs. A major challenge lies in the absence of a metric to quantify how well the pretrained generative ability is preserved. Current finetuning methods often implicitly presume that a smaller Euclidean distance between the finetuned and pretrained models corresponds to better retention of pretrained capabilities. For this reason, these methods usually employ very small learning rates. Although this assumption might hold in some cases, we contend that Euclidean distance alone cannot fully reflect the extent of semantic preservation. Consequently, adopting a more structural metric to assess differences between the finetuned and pretrained models could significantly enhance the retention of pretrained performance and improve finetuning stability.

Motivated by the empirical finding that hyperspherical similarity effectively captures semantic information~\cite{liu2017deep,liu2018decoupled,chen2020angular}, we employ hyperspherical energy~\cite{liu2018learning} to describe the pairwise relational structure among neurons. Hyperspherical energy is defined as the sum of hyperspherical similarity (such as cosine similarity) across all pairs of neurons within the same layer, reflecting the degree of neuron uniformity on the unit hypersphere~\cite{liu2021learning}. We hypothesize that a well-finetuned model should exhibit minimal change in hyperspherical energy relative to the pretrained model. A straightforward approach would be to introduce a regularizer to maintain hyperspherical energy during finetuning, but this does not guarantee effective minimization of the hyperspherical energy difference. Consequently, we leverage an invariance property of hyperspherical energy — the pairwise hyperspherical similarity remains unchanged when the same orthogonal transformation is applied to all neurons. Inspired by this invariance, we propose Orthogonal Finetuning (OFT), which adapts large text-to-image diffusion models to downstream tasks while preserving their hyperspherical energy. The key concept is to learn a layer-shared orthogonal transformation for neurons that maintains their pairwise angular relationships. OFT can also be interpreted as modifying the canonical coordinate system for neurons within the same layer. Additionally, considering that a smaller Euclidean distance between the finetuned and pretrained models indicates better retention of pretraining performance, we introduce a variant called Constrained Orthogonal Finetuning (COFT), which restricts the finetuned model to lie within a hypersphere of fixed radius centered at the pretrained neurons.

The rationale behind using orthogonal transformation for finetuning neurons partially stems from the 2D Fourier transform, where an image can be broken down into magnitude and phase spectra. The phase spectrum, representing the angular relationship between the input and the basis, retains most of the semantic content. For instance, an image's phase spectrum combined with a random magnitude spectrum can still reconstruct the original image without semantic loss. This observation implies that altering neuron directions is essential for semantically adjusting the generated image, which aligns with the objectives of both subject-driven and controllable generation. Nevertheless, permitting extensive changes in neuron directions can compromise the pretrained generative capabilities. To limit this flexibility, we propose preserving the angles between any pair of neurons, based on the premise that these inter-neuronal angles are vital for encoding the neural network’s knowledge. Following this insight, it is intuitive to learn an orthogonal transformation shared across neurons within each layer, ensuring the hyperspherical energy remains constant.

Additionally, we take inspiration from orthogonal over-parameterized training~\cite{liu2021orthogonal}, which trains classification networks from scratch by applying orthogonal transformations to a randomly initialized network. This approach is grounded in the fact that a randomly initialized network exhibits a provably low expected hyperspherical energy, and the goal in \cite{liu2021orthogonal} is to maintain low hyperspherical energy during training (since lower energy correlates with improved generalization in classification~\cite{liu2018learning,lin2020regularizing}). The work in \cite{liu2021orthogonal} demonstrates that orthogonal transformation provides sufficient flexibility to train neural networks that generalize well for classification tasks. In contrast, our focus is on finetuning text-to-image diffusion models to enhance controllability and downstream generative performance. We highlight two main differences between our OFT approach and \cite{liu2021orthogonal}. First, whereas \cite{liu2021orthogonal} aims to minimize hyperspherical energy, OFT seeks to preserve the pretrained model’s hyperspherical energy to maintain its intrinsic structure during finetuning. Minimizing hyperspherical energy in diffusion model finetuning risks destroying original semantic structures. Second, OFT attempts to minimize deviation from the pretrained model, leading to a constrained approach, while \cite{liu2021orthogonal} does not impose such restrictions. Achieving an optimal balance between flexibility and stability is crucial in finetuning, and we contend that our OFT framework effectively accomplishes this. Our contributions are summarized as follows:

\begin{itemize}[leftmargin=*,nosep]

    \item We introduce a new finetuning technique called Orthogonal Finetuning (OFT) aimed at enhancing the controllability of text-to-image diffusion models. To boost stability further, we present a constrained variant that restricts the angular deviation from the original pretrained model.
    \item Unlike current finetuning strategies, OFT adjusts model parameters while provably maintaining the hyperspherical energy, which we empirically observe as a crucial indicator of preserving the generative semantics of the pretrained model.
    \item We deploy OFT in two scenarios: subject-driven generation and controllable generation. Through extensive empirical evaluation, we show substantial improvements over previous methods regarding generation quality, convergence speed, and finetuning stability. Additionally, OFT demonstrates greater sample efficiency, converging effectively with significantly fewer training images and epochs.
    \item For controllable generation, we propose a novel control consistency metric to assess controllability. The key concept is to derive the control signal from the generated image and then compare it with the original control input. This metric further confirms OFT’s strong controllability.
\end{itemize}

\section{Related Work}

\textbf{Text-to-image diffusion models}. Significant advancements~\cite{saharia2022photorealistic,ramesh2022hierarchical,rombach2022high,gu2022vector,nichol2022glide} have been achieved in text-to-image synthesis, largely propelled by the rapid evolution of diffusion-based generative models~\cite{ho2020denoising,song2021score,song2021denoising,dhariwal2021diffusion} alongside vision-language representation learning~\cite{su2020vlbert,lu2019vilbert,radford2021learning,wang2021simvlm,li2022blip,li2023blip,alayrac2022flamingo,schuhmann2022laion}. GLIDE~\cite{nichol2022glide} and Imagen~\cite{saharia2022photorealistic} perform diffusion modeling directly in pixel space, where GLIDE trains the text encoder jointly with a diffusion prior using paired text-image data, whereas Imagen employs a frozen pretrained text encoder. Stable Diffusion~\cite{rombach2022high} and DALL-E2~\cite{ramesh2022hierarchical} conduct diffusion modeling within latent spaces. Stable Diffusion utilizes VQ-GAN~\cite{esser2021taming} to learn a visual codebook as the latent space, while DALL-E2 leverages CLIP~\cite{radford2021learning} to create a shared latent embedding for images and text. Beyond diffusion models, generative adversarial networks~\cite{reed2016generative,zhang2017stackgan,xu2018attngan,li2019controllable} and autoregressive models~\cite{ramesh2021zero,ding2021cogview,wu2022nuwa,yuscaling} have also been explored for text-to-image generation. OFT is intrinsically a model-agnostic finetuning approach and can be utilized with any text-to-image diffusion model.

\textbf{Subject-driven generation}. To avoid altering the subject, \cite{avrahami2022blended,nichol2022glide} incorporate user-provided masks as additional conditioning. Inversion methods~\cite{choi2021ilvr,dhariwal2021diffusion,ramesh2022hierarchical,gal2022image} allow editing the context while keeping the subject unchanged. \cite{hertz2022prompt} achieves local and global edits without needing input masks. However, these methods struggle to retain identity-specific details of the subject. Pivotal Tuning~\cite{roich2022pivotal} finetunes the generator around an initial inverted latent code with added regularization to maintain identity. Similarly, \cite{nitzan2022mystyle} trains a personalized generative face prior from a set of images of an individual. \cite{casanova2021instance} can produce various instance-specific outputs but may sacrifice detailed fidelity. Using a customized token and a few subject images, DreamBooth~\cite{ruiz2023dreambooth} finetunes a text-to-image diffusion model by combining a reconstruction loss with a class-specific prior preservation loss. OFT builds upon the DreamBooth framework but replaces naive low-rate finetuning with orthogonal transformations during the finetuning process.

\textbf{Controllable generation}. Image-to-image translation can be regarded as a type of controllable generation, with earlier approaches primarily utilizing conditional generative adversarial networks~\cite{isola2017image,zhu2017toward,wang2018high,park2019semantic,choi2018stargan}. Diffusion models have also been applied for image-to-image translation~\cite{saharia2022palette,voynov2022sketch,wang2022pretraining}. More recently, ControlNet~\cite{zhang2023adding} introduces a method to guide a pretrained diffusion model by fine-tuning it to incorporate additional control signals, achieving notable performance in controllable generation. A similar contemporary work, T2I-Adapter~\cite{mou2023t2i}, also fine-tunes pretrained diffusion models to enhance control over generated images. Aligning with the task settings in \cite{zhang2023adding,mou2023t2i}, we utilize OFT to fine-tune pretrained diffusion models, consistently improving controllability while requiring less training data and fewer fine-tuning parameters. Importantly, OFT does not impose any extra computational cost during inference.

\textbf{Model finetuning}. The practice of fine-tuning large pretrained models for downstream applications has grown increasingly prevalent~\cite{devlin2018bert,brown2020language,he2022masked}. Adaptation techniques, such as those in \cite{hulora2022,houlsby2019parameter,pfeiffer2020adapterfusion}, represent a form of fine-tuning and have been extensively researched in natural language processing. LoRA~\cite{hulora2022} is the work most closely related to OFT, as it assumes a low-rank structure in the additive weight updates during fine-tuning. By contrast, OFT employs a layer-shared orthogonal transformation that multiplicatively adjusts neuron weights, and it provably maintains the pairwise angles among neurons within the same layer, leading to enhanced stability.

\section{Orthogonal Finetuning}

\subsection{Why Does Orthogonal Transformation Make Sense?}

We begin by exploring why orthogonal transformation is advantageous for fine-tuning text-to-image diffusion models. This question is broken down into two parts: (1) the rationale for fine-tuning the angles of neurons (i.e., their directions), and (2) the reason for using orthogonal transformations to adjust these angles.

Regarding the first question, we take inspiration from the empirical findings in \cite{liu2018decoupled,chen2020angular} which indicate that angular differences in features effectively capture the semantic gap. SphereNet~\cite{liu2017deep} demonstrates that training a neural network with all neurons normalized on a unit hypersphere can achieve comparable capacity along with improved generalizability, suggesting that the direction of neurons encapsulates the most critical information from the data. To further illustrate the significance of neuron angles, we carry out a simple experiment depicted in Figure~\ref{fig:toy_ae}, where we train a conventional convolutional autoencoder on a set of flower images. During training, the standard inner product is used to produce the feature map ($z$ represents the element output from the convolution kernel $\bm{w}$, and $\bm{x}$ is the input within the sliding window). In testing, we evaluate three different approaches to generate the feature map: (a) the inner product as used in training, (b) only the magnitude information, and (c) solely the angular information. The outcomes shown in Figure~\ref{fig:toy_ae} reveal that angular information from neurons can almost perfectly reconstruct the input images, whereas the magnitude contains no meaningful information. It is important to highlight that we do not apply the cosine activation (c) during training; training is performed solely with the inner product. These findings suggest that neuron angles (directions) predominantly store the semantic content of input images. Consequently, fine-tuning neuron directions may be more effective for altering the semantic information of images.

Addressing the second question, the most straightforward method to fine-tune neuron directions is to simultaneously rotate or reflect all neurons within the same layer, naturally leading to orthogonal transformations. While employing other angular transformations that rotate different neurons by varying angles could offer more flexibility, we observe that orthogonal transformations provide an optimal balance between flexibility and structural regularity. Additionally, \cite{liu2021orthogonal} establishes that orthogonal transformations possess sufficient power for learning neural networks. To reinforce our point, we conduct an experiment to showcase the regularization benefits induced by enforcing orthogonality. Using the controllable generation setup from ControlNet~\cite{zhang2023adding}, the results are presented in Figure~\ref{fig:ortho}. Our standard OFT method demonstrates stability and achieves precise control once training concludes (epoch 20). In contrast, OFT without the orthogonality constraint fails to generate any realistic images and does not achieve control. This experiment confirms the crucial role of the orthogonality constraint in OFT.

\subsection{General Framework}

The traditional finetuning method usually employs gradient descent with a reduced learning rate to update a model (or specific layers of it). This small learning rate implicitly encourages the finetuned model to remain close to the pretrained model, essentially aiming to train the model by implicitly minimizing $\thickmuskip=2mu \medmuskip=2mu \| \bm{M} - \bm{M}^0\|$, where $\bm{M}$ represents the finetuned model parameters and $\bm{M}^0$ denotes the pretrained weights. While this implicit constraint is intuitively reasonable, it can still allow too much flexibility when finetuning a large model. To tackle this, LoRA imposes an additional low-rank constraint on the weight updates, namely $\thickmuskip=2mu \medmuskip=2mu \text{rank}(\bm{M} - \bm{M}^0)=r'$, with $r'$ chosen to be small. Unlike LoRA, OFT enforces a constraint on the pairwise neuron similarity by requiring $\thickmuskip=2mu \medmuskip=2mu \| \text{HE}(\bm{M})-\text{HE}(\bm{M}^0) \|=0$, where $\text{HE}(\cdot)$ stands for the hyperspherical energy of the weight matrix. 

To illustrate, consider a fully connected layer $\thickmuskip=2mu \medmuskip=2mu \bm{W}=\{\bm{w}_1,\cdots,\bm{w}_n\}\in\mathbb{R}^{d\times n}$, where each $\bm{w}_i\in\mathbb{R}^d$ is the $i$-th neuron ($\bm{W}^0$ denotes the pretrained weights). The output vector $\thickmuskip=2mu \medmuskip=2mu \bm{z}\in\mathbb{R}^n$ of this layer is computed as $\thickmuskip=2mu \medmuskip=2mu \bm{z}=\bm{W}^\top \bm{x}$, where $\thickmuskip=2mu \medmuskip=2mu \bm{x} \in \mathbb{R}^d$ is the input. OFT can be understood as minimizing the difference in hyperspherical energy between the finetuned and pretrained models:

\begin{equation}\label{eq:oft_principle}
\footnotesize
\min_{\bm{W}} \left\| \text{HE}(\bm{W})-\text{HE}(\bm{W}^0)\right\|~~\Leftrightarrow~~\min_{\bm{W}} \bigg{\|} \sum_{i\neq j} \|\hat{\bm{w}}_i-\hat{\bm{w}}_j\|^{-1}-\sum_{i\neq j} \|\hat{\bm{w}}^0_i-\hat{\bm{w}}^0_j\|^{-1}\bigg{\|}
\end{equation}

Here, $\thickmuskip=2mu \medmuskip=2mu \hat{\bm{w}}_i := \bm{w}_i/\|\bm{w}_i\|$ denotes the normalized $i$-th neuron, and the hyperspherical energy for a fully connected layer $\bm{W}$ is defined as $\thickmuskip=2mu \medmuskip=2mu \text{HE}(\bm{W}):= \sum_{i\neq j} \|\hat{\bm{w}}_i-\hat{\bm{w}}_j\|^{-1}$. It is straightforward to see that the minimum attainable value in Eq.~\eqref{eq:oft_principle} is zero. This minimum is achieved if and only if $\bm{W}$ differs from $\bm{W}^0$ by a rotation or reflection, i.e., $\thickmuskip=2mu \medmuskip=2mu \bm{W} = \bm{R}\bm{W}^0$, where $\thickmuskip=2mu \medmuskip=2mu \bm{R}\in\mathbb{R}^{d \times d}$ is an orthogonal matrix (with determinant $1$ or $-1$ corresponding to rotation or reflection, respectively). This concept underlies OFT: finetuning the network by learning orthogonal matrices shared across layers that transform neurons within each layer.

\begin{equation}\label{eq:oft_general}
    \footnotesize
    \bm{z}=\bm{W}^\top\bm{x}=(\bm{R}\cdot\bm{W}^0)^\top\bm{x},~~~\text{s.t.}~\bm{R}^\top\bm{R}=\bm{R}\bm{R}^\top=\bm{I}
\end{equation}

Here, $\bm{W}$ represents the weight matrix fine-tuned by OFT, and $\bm{I}$ is the identity matrix. The OFT method is depicted in Figure~\ref{fig:block}. In a manner akin to zero initialization in LoRA, it is essential to ensure that OFT fine-tunes the pretrained model starting exactly from $\bm{W}^0$. To this end, the orthogonal matrix $\bm{R}$ is initialized as the identity matrix, ensuring the fine-tuned model begins with the pretrained weights. To maintain the orthogonality of $\bm{R}$, differential orthogonalization techniques as outlined in \cite{liu2021orthogonal,lezcano2019cheap} can be employed. We will further elaborate on how to enforce orthogonality in an efficient computational manner.

\subsection{Efficient Orthogonal Parameterization}\label{sect:eff_ortho}

Although the Gram-Schmidt orthogonalization method is differentiable, it tends to be computationally costly in practice~\cite{liu2021orthogonal}. To improve efficiency, we utilize the Cayley parameterization to generate the orthogonal matrix. Concretely, the orthogonal matrix is formed as $\thickmuskip=2mu \medmuskip=2mu \bm{R}=(\bm{I}+\bm{Q})(\bm{I}-\bm{Q})^{-1}$ where $\bm{Q}$ is a skew-symmetric matrix fulfilling $\thickmuskip=2mu \medmuskip=2mu \bm{Q}=-\bm{Q}^\top$. This approach gains efficiency at a minor cost — the Cayley parameterization produces only orthogonal matrices with determinant $1$, belonging to the special orthogonal group. Fortunately, this restriction does not adversely impact performance in practice. Even with the Cayley transform for parameterizing $\bm{R}$, when the dimension $d$ is large, $\bm{R}$ can still be parameter-heavy. To mitigate this, we propose representing $\bm{R}$ as a block-diagonal matrix consisting of $r$ blocks, yielding the structure:

\begin{equation}
    \footnotesize
    \bm{R}=\text{diag}(\bm{R}_1,\bm{R}_2,\cdots,\bm{R}_r)=\begin{bmatrix}
    \bm{R}_1\in \text{O}(\frac{d}{r}) &  &\\
    &\ddots & \\
    & & \bm{R}_r\in \text{O}(\frac{d}{r})
    \end{bmatrix}\in \text{O}(d)
\end{equation}

where $\text{O}(d)$ represents the orthogonal group in dimension $d$, and $\thickmuskip=2mu \medmuskip=2mu \bm{R}\in\mathbb{R}^{d\times d}$ along with $\thickmuskip=2mu \medmuskip=2mu \bm{R}_i\in\mathbb{R}^{d/r\times d/r},\forall i$ are orthogonal matrices. When $\thickmuskip=2mu \medmuskip=2mu r=1$, the block-diagonal orthogonal matrix reduces to a conventional unconstrained orthogonal matrix. For an orthogonal matrix sized $\thickmuskip=2mu \medmuskip=2mu d\times d$, the parameter count is $\thickmuskip=2mu \medmuskip=2mu d(d-1)/2$, yielding a computational complexity of $\mathcal{O}(d^2)$. In the case of an $r$-block diagonal orthogonal matrix, the parameter count becomes $\thickmuskip=2mu \medmuskip=2mu d(d/r-1)/2$, which corresponds to a complexity of $\thickmuskip=2mu \medmuskip=2mu \mathcal{O}(d^2/r)$. Additionally, one may share the block matrices to further lower the number of parameters, i.e., $\thickmuskip=2mu \medmuskip=2mu \bm{R}_i=\bm{R}_j,\forall i\neq j$. This reduces the parameter complexity to $\mathcal{O}(d^2/r^2)$. Despite these parameter-efficiency improvements, the resultant matrix $\bm{R}$ remains orthogonal, ensuring no compromise in preserving hyperspherical energy.

We compare the parameter efficiency of OFT with LoRA. For LoRA with low-rank parameter $r'$, the count of trainable parameters is $\thickmuskip=2mu \medmuskip=2mu r'(d+n)$. If both $r$ and $r'$ scale with neuron dimension $d$ (e.g., $\thickmuskip=2mu \medmuskip=2mu r=r'=\alpha d$ for some constant $0<\alpha\leq 1$), then LoRA’s parameter complexity is $\mathcal{O}(d^2 + dn)$, whereas OFT’s complexity is $\mathcal{O}(d)$. To illustrate the difference concretely, consider a weight matrix of size $\thickmuskip=2mu \medmuskip=2mu 128 \times 128$: LoRA has $2,048$ trainable parameters with $\thickmuskip=2mu \medmuskip=2mu r'=8$, while OFT has $960$ trainable parameters with $\thickmuskip=2mu \medmuskip=2mu r=8$ (without block sharing).

\subsection{Constrained Orthogonal Finetuning}

We can further restrict the flexibility of the original OFT by confining the finetuned model to remain close to the pretrained model. Specifically, COFT employs the following forward pass:

\begin{equation}\label{eq:coft_general}
    \footnotesize
    \bm{z} = \bm{W}^\top \bm{x} = (\bm{R} \cdot \bm{W}^0)^\top \bm{x}, \quad \text{s.t.} \quad \bm{R}^\top \bm{R} = \bm{R} \bm{R}^\top = \bm{I}, \quad \norm{\bm{R} - \bm{I}} \leq \epsilon
\end{equation}

which enforces an orthogonality constraint alongside an $\epsilon$-bounded deviation from the identity matrix. The orthogonality constraint can be fulfilled using the Cayley parameterization introduced in Section~\ref{sect:eff_ortho}. However, integrating the $\epsilon$-deviation constraint into the Cayley-parameterized orthogonal matrix is not straightforward. To gain further insight into the Cayley transform, we use the Neumann series to approximate $\thickmuskip=2mu \medmuskip=2mu \bm{R} = (\bm{I} + \bm{Q})(\bm{I} - \bm{Q})^{-1}$ by $\thickmuskip=2mu \medmuskip=2mu \bm{R} \approx \bm{I} +

series converges under the operator norm). Hence, we are able to incorporate the constraint $\thickmuskip=2mu \medmuskip=2mu\|\bm{R}-\bm{I}\|\leq\epsilon$ within the Cayley transform, resulting in an equivalent constraint $\thickmuskip=2mu \medmuskip=2mu \|\bm{Q}-\bm{0}\|\leq\epsilon'$, where $\bm{0}$ represents the zero matrix and $\epsilon'$ is a distinct error hyperparameter (different from $\epsilon$). This newly imposed restriction on the matrix $\bm{Q}$ can be straightforwardly enforced via projected gradient descent. To initialize the orthogonal matrix $\bm{R}$ as the identity, we set $\bm{Q}$ initially as the zero matrix. COFT can be interpreted as the combination of two explicit constraints: minimizing hyperspherical energy difference and limiting deviation from the pretrained model. The latter constraint is typically implicit in current finetuning approaches, but COFT explicitly incorporates it. Although OFT demonstrates strong performance, we find that COFT provides enhanced finetuning stability compared to OFT, attributable to this explicit deviation constraint. Figure~\ref{fig:eps_coft} illustrates an example showing how the parameter $\epsilon$ influences COFT's performance. It is evident that $\epsilon$ regulates finetuning flexibility: a larger $\epsilon$ results in a COFT-finetuned model closer to the OFT-finetuned model, while a smaller $\epsilon$ causes the COFT-finetuned model to behave increasingly like the original pretrained text-to-image diffusion model.

\subsection{Re-scaled Orthogonal Finetuning}

We introduce a straightforward extension of the original OFT by incorporating a learnable magnitude scaling factor for each neuron. This extension is motivated by the observation that rescaling neurons does not affect the hyperspherical energy, since magnitude normalization removes this influence. Concretely, the forward pass is defined as $\thickmuskip=2mu \medmuskip=2mu\bm{z}=(\bm{R}\bm{W}^0\bm{D})^\top\bm{x}$\footnote{\textbf{Errata}: In the NeurIPS camera-ready version, the forward pass for re-scaled OFT was erroneously stated as $\thickmuskip=2mu \medmuskip=2mu\bm{z}=(\bm{D}\bm{R}\bm{W}^0)^\top\bm{x}$. The original implementation remains correct; therefore, the results in Appendix~\ref{app:rescaled} remain valid.} where $\thickmuskip=2mu \medmuskip=2mu \bm{D}=\text{diag}(s_1,\cdots,s_n)\in\mathbb{R}^{n\times n}$ is a learnable diagonal matrix with strictly positive diagonal elements $s_1,\cdots,s_n$. Unlike the original OFT forward pass in Eq.~\eqref{eq:oft_general}, where only $\bm{R}$ is learned, both the diagonal matrix $\bm{D}$ and the orthogonal matrix $\bm{R}$ are now learned parameters. This re-scaled OFT enhances OFT's flexibility with a negligible increase in parameter count. Our experiments primarily use the original OFT to demonstrate the efficacy of orthogonal transformations alone, but we find that re-scaled OFT generally performs better (see Appendix~\ref{app:rescaled}).

\section{Intriguing Insights and Discussions}

\textbf{OFT is architecture-agnostic}. In principle, OFT can be applied to any neural network architecture. For Transformers, LoRA commonly targets attention weights~\cite{hulora2022}. To ensure a fair comparison with LoRA, we limit OFT finetuning to attention weights in our experiments. Beyond fully connected layers, OFT is also well suited for finetuning convolutional layers because the block-diagonal structure of $\bm{R}$ admits meaningful interpretations in convolutional contexts (unlike LoRA). When the number of blocks matches the number of input channels, each block transforms a single neuron channel independently, resembling depth-wise convolution kernels~\cite{chollet2017xception}. If all blocks in $\bm{R}$ are shared, OFT applies an orthogonal transformation across channels using a common matrix. We present a preliminary investigation on finetuning convolutional layers with OFT in Appendix~\ref{app:convolution}.

\textbf{Connection to LoRA}. LoRA adds a low-rank matrix to prevent significant shifts in the information contained within the pretrained weight matrix. In contrast, OFT enforces that the transformation applied to the pretrained weight matrix is orthogonal (full-rank), which safeguards the pretrained information from being disrupted. We can express OFT's forward pass as $\thickmuskip=2mu \medmuskip=2mu \bm{z}=(\bm{R}\bm{W}^0)^\top\bm{x}=(\bm{W}^0+(\bm{R}-\bm{I})\bm{W}^0)^\top\bm{x}$, where $\thickmuskip=2mu \medmuskip=2mu (\bm{R}-\bm{I})\bm{W}^0$ plays a role analogous to LoRA's low-rank weight update. Because $\bm{W}^0$ is generally full-rank, OFT similarly performs a low-rank weight update whenever $\thickmuskip=2mu \medmuskip=2mu \bm{R}-\bm{I}$ is low-rank. Just as LoRA employs a rank parameter $r'$, OFT uses a diagonal block parameter $r$ to reduce the count of trainable parameters. More intriguingly, LoRA and OFT represent two different strategies for parameter efficiency: LoRA leverages low-rank structure to minimize trainable parameters, whereas OFT adopts a distinct approach by utilizing sparsity structure, specifically block-diagonal orthogonality.

\textbf{Why OFT converges faster}. From one perspective, Figure~\ref{fig:toy_ae} shows that the most effective way to alter semantics is by modifying neuron directions, which aligns perfectly with the design of OFT. Additionally, OFT can be interpreted as fine-tuning neurons on a smooth hyperspherical manifold, leading to a more favorable optimization landscape. This observation is also supported empirically in \cite{liu2021orthogonal}.

\textbf{Why not minimize hyperspherical energy}. A crucial distinction from \cite{liu2021orthogonal} is that we do not strive to minimize hyperspherical energy. For classification tasks, it is preferable for neurons to avoid redundancy. Minimizing hyperspherical energy arranges all neurons uniformly across the hypersphere, which is not a suitable objective for finetuning as it risks destroying the pretrained information.

\textbf{Balancing flexibility and regularity in finetuning}. We identify a fundamental trade-off between flexibility and regularity. Traditional finetuning offers the greatest flexibility but tends to have poor stability and is prone to model collapse. Surprisingly, OFT, with its simplicity, strikes a good compromise by maintaining pairwise neuron angles. Additionally, the block-diagonal parameterization can be interpreted as a stronger form of regularization on the orthogonal matrix.

\textbf{No extra inference cost}. Unlike ControlNet, the OFT framework does not add any inference overhead to the finetuned model. During inference, we can directly multiply the learned orthogonal matrix $\bm{R}$ with the pretrained weight matrix $\bm{W}^0$ to produce an equivalent weight matrix $\thickmuskip=2mu \medmuskip=2mu \bm{W}=\bm{R}\bm{W}^0$. Consequently, the inference speed remains identical to that of the pretrained model.

\section{Experiments and Results}

\textbf{General setup}. For our experiments, we use Stable Diffusion v1.5~\cite{rombach2022high} as the pretrained text-to-image model. To ensure fairness, we randomly select generated images from each method. In subject-driven generation, we largely follow the approach of DreamBooth~\cite{ruiz2023dreambooth}. For controllable generation, we adhere mostly to ControlNet~\cite{zhang2023adding} and T2I-Adapter~\cite{mou2023t2i}. To maintain a fair comparison with LoRA, OFT or COFT is applied only to the identical layer where LoRA is used. Additional results and details are provided in Appendix~\ref{app:settings}.

\subsection{Subject-driven Generation}

\textbf{Setup}. DreamBooth~\cite{ruiz2023dreambooth} and LoRA~\cite{hulora2022} serve as our baseline methods. All approaches utilize the same loss function as DreamBooth. For DreamBooth and LoRA, we generally follow the hyperparameter settings recommended in their original papers. Further results can be found in Appendix~\ref{app:settings}, \ref{app:coft_oft}, \ref{app:more_qualitative}, \ref{app:failure}.

\textbf{Stability and convergence during finetuning}. We begin by assessing the finetuning stability and convergence rate for DreamBooth, LoRA, OFT, and COFT. The outcomes are illustrated in Figure~\ref{fig:teaser} and Figure~\ref{fig:db_convergence}. It is evident that both COFT and OFT can finetune the diffusion model with considerable stability. After 400 iterations, both DreamBooth and OFT variants exhibit effective control, whereas LoRA fails to maintain the subject's identity. After 2000 iterations, DreamBooth begins to produce collapsed images, and LoRA is unable to generate the yellow shirt, producing yellow fur instead. Conversely, OFT and COFT continue to maintain stable and consistent control over the generated images. These findings confirm the rapid yet reliable convergence of our OFT framework in subject-driven generation. We emphasize that the robustness to the number of finetuning iterations is crucial, as it helps mitigate the difficulty of tuning iteration counts for various subjects. For both OFT and COFT, it is possible to directly select a relatively large number of iterations without meticulous tuning. In the case of COFT with an appropriate $\epsilon$, setting both the learning rate and iteration count becomes straightforward.

\textbf{Quantitative evaluation}. Following \cite{ruiz2023dreambooth}, we perform a quantitative assessment of subject fidelity (DINO~\cite{caron2021emerging}, CLIP-I~\cite{radford2021learning}), text prompt fidelity (CLIP-T~\cite{radford2021learning}), and sample diversity (LPIPS~\cite{zhang2018unreasonable}). CLIP-I measures the average pairwise cosine similarity of CLIP embeddings between generated and real images. DINO is analogous to CLIP-I but uses ViT S/16 DINO embeddings instead. CLIP-T calculates the average cosine similarity between the text prompt and generated images' CLIP embeddings. We also assess the average LPIPS cosine similarity among generated images of the same subject with identical text prompts. Table~\ref{table:dreambooth_final} demonstrates that both COFT and OFT surpass DreamBooth and LoRA by a significant margin in the DINO and CLIP-I metrics, while achieving slightly better or comparable results in prompt fidelity and diversity metrics. To reduce randomness, each method is finetuned repeatedly with 30 different random seeds on the same pretrained model. The results indicate that our OFT framework not only ensures better convergence and stability but also consistently produces superior final outcomes.

\textbf{Qualitative comparison}. To better illustrate the advantages of OFT, we present several randomly selected examples for subject-driven generation in Figure~\ref{fig:dreambooth_final}. For an unbiased comparison, all examples are generated using the same finetuned model corresponding to each method, ensuring that no individual text prompt is optimized specifically for the final outputs. For each approach, we choose the model that achieves the highest validation CLIP scores. As shown in Figure~\ref{fig:dreambooth_final}, both OFT and COFT demonstrate strong semantic consistency in preserving the subject, whereas LoRA frequently struggles to maintain the subject’s identity (for instance, LoRA completely loses the subject identity in the bowl example). Additionally, OFT and COFT exhibit much finer control via text prompts, while DreamBooth, although good at preserving subject identity, often fails to produce images that align well with the textual prompt (e.g., the first row of the bowl example). This qualitative evaluation highlights that our OFT framework offers superior controllability and subject preservation simultaneously. Furthermore, OFT is robust to the number of iterations, maintaining performance even with many iterations, unlike DreamBooth and LoRA. Additional qualitative examples are provided in Appendix~\ref{app:more_qualitative}. A human evaluation conducted in Appendix~\ref{app:human} further supports the advantage of OFT.

\subsection{Controllable Generation}

\textbf{Settings}. We compare against ControlNet~\cite{zhang2023adding}, T2I-Adapter~\cite{mou2023t2i}, and LoRA~\cite{hulora2022} as baseline methods. We focus on three challenging controllable generation tasks in the main paper: Canny edge to image (C2I) on the COCO dataset~\cite{lin2014microsoft}, segmentation map to image (S2I) on the ADE20K dataset~\cite{zhou2017scene}, and landmark to face (L2F) on the CelebA-HQ dataset~\cite{karras2018progressive,xia2021tedigan}. Each method fine-tunes Stable Diffusion (SD) v1.5 on these three datasets for 20 epochs. Further results are available in Appendices~\ref{app:more_qualitative}, \ref{app:more_control_task}, and \ref{app:failure}.

\textbf{Convergence}. We assess the convergence rate of ControlNet, T2I-Adapter, LoRA, and COFT on the L2F task through both quantitative and qualitative analyses. For the quantitative evaluation, we calculate the average $\ell_2$ distance between the control face landmarks and the predicted face landmarks. In Figure~\ref{fig:face_convergence}, we display the face landmark error of models fine-tuned over various numbers of epochs. It is evident that both COFT and OFT converge significantly faster. While LoRA requires 20 epochs to reach the performance level of our OFT framework at the 8th epoch, OFT and COFT achieve this much earlier. It is noteworthy that OFT and COFT utilize a comparable amount of trainable parameters as LoRA (substantially fewer than ControlNet), yet they converge much more efficiently than current methods. Additionally, the rapid convergence of OFT is further supported by the results in Figure~\ref{fig:teaser}. The rightmost example there demonstrates that OFT is considerably more data-efficient compared to ControlNet and LoRA, as it can effectively converge with just 5\% of the ADE20K dataset. For qualitative comparisons, we concentrate on OFT, COFT, and ControlNet, since ControlNet achieves landmark errors closest to ours. The results presented in Figure~\ref{fig:face_convergence_qual} indicate that both OFT and COFT converge steadily, with the generated face pose progressively aligning to the control landmarks. In contrast, the controllability in ControlNet abruptly appears after the 8th epoch, which aligns with the sudden convergence pattern described in \cite{zhang2023adding}. This stability in convergence makes the OFT framework much easier to apply in practice, as its training progression is smoother and more predictable. Consequently, identifying suitable hyperparameters for OFT becomes simpler.

\textbf{Quantitative comparison}. We propose a control consistency metric to assess the effectiveness of controllable generation. The core concept involves extracting the control signal from the generated image and then comparing it against the original input control signal. For the C2I task, we calculate IoU and F1 score. For the S2I task, we measure mean IoU, as well as mean and overall accuracy. For the L2F task, we evaluate the average $\ell_2$ distance between the control landmarks and the predicted landmarks. Further details on the consistency metrics can be found in Appendix~\ref{app:settings}. For all methods compared, we apply the best possible hyperparameter settings. The results presented in Table~\ref{table:control_gen} demonstrate that both OFT and COFT achieve significantly stronger and more precise control than the other techniques. We note that adapter-based methods (\emph{e.g.}, T2I-Adapter and ControlNet) show slower convergence and produce inferior final outcomes. Compared to ControlNet, LoRA performs better on the S2I task but worse on the C2I and L2F tasks. Overall, we observe that LoRA’s performance ceiling is relatively low, even with careful hyperparameter tuning. In contrast, the performance of our OFT framework has not yet plateaued, as we empirically observe continued improvement with an increasing number of trainable parameters. We highlight that our quantitative evaluation of controllable generation represents a novel contribution, as it enables precise assessment of control performance for finetuned models (subject to the accuracy of the off-the-shelf segmentation/detection models).

\textbf{Qualitative comparison}. We further conduct a qualitative comparison of OFT and COFT against the current leading methods, such as ControlNet, T2I-Adapter, and LoRA. The randomly generated images shown in Figure~\ref{fig:control_final} illustrate that OFT and COFT not only produce high-fidelity and realistic images but also provide precise control. In the S2I task, it is evident that LoRA fails to generate images consistent with the input segmentation map, whereas ControlNet, OFT, and COFT successfully control the generated images. Compared to ControlNet, both OFT and COFT generate high-quality images featuring richer details and more accurate geometric structures, despite using significantly fewer model parameters. In the C2I task, OFT and COFT can hallucinate semantically relevant images from coarse Canny edges, whereas T2I-Adapter and LoRA perform notably worse. For the L2F task, our approach achieves the most precise pose control on generated faces, even with challenging poses. Across all three control tasks, we demonstrate that OFT and COFT produce qualitatively superior images compared to the leading baseline methods, confirming the effectiveness of the OFT framework in controllable image generation. To provide a more thorough qualitative comparison, additional examples for all three control tasks are included in Appendix~\ref{app:control_exp}. Furthermore, we show that OFT performs well on additional control tasks (such as dense pose to human body, sketch to image, and depth to image) in Appendix~\ref{app:more_control_task}.

\section{Concluding Remarks and Open Problems}

Inspired by the insight that angular relationships between neurons play a key role in defining visual semantics, we introduce a straightforward yet powerful finetuning technique—orthogonal finetuning—for controlling text-to-image diffusion models. Our approach specifically focuses on two text-to-image tasks: subject-driven generation and controllable generation. In comparison to existing approaches, OFT offers enhanced controllability and greater finetuning stability while using fewer finetuning parameters. Crucially, OFT does not add any extra inference cost, making the resulting model efficient and easy to deploy.

OFT also raises several intriguing open questions. First, the enforcement of orthogonality in OFT relies on Cayley parametrization, which requires computing a matrix inverse. This step somewhat restricts the scalability of OFT. Although we partially mitigate this limitation by employing a block diagonal parametrization, accelerating the matrix inversion in a differentiable manner remains a challenging problem. Second, OFT holds promising potential for compositionality, given that the orthogonal matrices generated by different OFT finetuning tasks can be multiplied together and still yield an orthogonal matrix. Investigating whether this set of orthogonal matrices preserves the knowledge from all the downstream tasks is an interesting avenue for future research. Lastly, the parameter efficiency of OFT heavily depends on the block diagonal structure, which inevitably introduces certain biases and restricts flexibility. Developing methods to enhance parameter efficiency in a more effective and less biased manner remains a crucial open challenge.

\end{document}